% File: abstract_ru.tex
% ReconX diploma --- Russian annotation (Аннотация)
\selectlanguage{russian}
\chapter*{Аннотация}
\addcontentsline{toc}{chapter}{Аннотация}
\markboth{\foreignlanguage{russian}{АННОТАЦИЯ}}{\foreignlanguage{russian}{АННОТАЦИЯ}}
\begin{SingleSpace}
В данном дипломном проекте представлены проектирование, реализация и экспериментальная оценка \textbf{ReconX} --- фреймворка для внешнего тестирования веб-приложений на проникновение. Предложенный фреймворк следует методологии <<чёрного ящика>> и объединяет пассивную разведку, активное снятие отпечатков, анализ конфигурации и широкий спектр модулей проверки уязвимостей в отношении удалённой цели, исходный код и учётные данные которой недоступны тестировщику.

Фреймворк реализован на \textbf{Python~3} и организован в виде конвейера из 46 специализированных модулей, сгруппированных в 15 этапов исполнения. Модули разведки собирают информацию о поддоменах, DNS-записях, портах и используемых технологиях; конфигурационные модули проводят аудит \textbf{TLS}, заголовков безопасности, \textbf{CORS}, аутентификационных cookie-файлов и \textbf{JWT}-токенов; активные модули проверяют наиболее актуальные категории \textbf{OWASP Top~10}, включая отражённый XSS, SQL-инъекции, IDOR, инъекции заголовка Host, отравление кэша, загрязнение прототипа, открытые редиректы, SSRF/SSTI, XXE, HTTP request smuggling и состояния гонки. Центральный \textit{реестр находок} нормализует каждый результат в единую схему с метаданными \textbf{CWE}, \textbf{CVSS}, OWASP и MITRE ATT\&CK; коррелятор и модуль рисков активов затем приоритизируют наиболее уязвимые активы. Локальная \textbf{большая языковая модель} (Ollama / DeepSeek-R1) формирует пояснительный аналитический отчёт по завершении сканирования.

Экспериментальная оценка была проведена в двух взаимодополняющих условиях: контролируемый лабораторный запуск против заведомо уязвимого \textit{OWASP Juice Shop} и санкционированное реальное исследование академической системы \texttt{miras.app}, выполненное в режиме \texttt{bug\_bounty} фреймворка (\texttt{allow\_write: false}). Фреймворк завершил полный конвейер на реальной цели приблизительно за 1\,мин~43\,с и сформировал 721 уникальную находку по категориям OWASP Top~10, ранжировал 49 обнаруженных активов на четыре уровня риска (включая 5 активов \textit{критического} уровня) и сгенерировал две сводные межмодульные находки через коррелятор. Обнаруженные уязвимости соответствовали ожидаемым категориям обеих тестовых целей, лишь 1 из 4\,227 отправленных HTTP-запросов оказался за пределами настроенной области (и был заблокирован фреймворком до отправки), а сгенерированная ИИ исполнительная сводка корректно вывела актив с наивысшим риском в начало отчёта.

Результаты подтверждают, что предложенный метод делает внешнее тестирование веб-приложений на проникновение более воспроизводимым, более объяснимым и существенно более эффективным по времени по сравнению с полностью ручным подходом, оставляя за оператором контроль над интрузивными действиями через явную систему профилей (\textit{safe}, \textit{bug\_bounty}, \textit{deep}, \textit{intrusive}).

\par\noindent{\small\textbf{Ключевые слова:} внешнее тестирование на проникновение; безопасность веб-приложений; OWASP Top~10; разведка; сканирование уязвимостей; корреляция находок; анализ безопасности с помощью LLM; ReconX.}
\end{SingleSpace}
\selectlanguage{english}
\clearpage
